\section{Paying for the elementary projection on the selected scales}
\label{sec:projection-payment-producer}
We now verify the principal reserve and the numerical projection
allowance used in Section~\ref{sec:leaf-transfer}. All quantities
refer to the marked expansion constructed in
Sections~\ref{sec:marked-block-integral} and
\ref{sec:marked-activity-jets}. In particular the regular logarithmic
terms and the differentiated forced array are those of the same
background transition. We retain the incoming sector and path
estimates of that construction.

\subsection{The singleton identity before evaluating the gas}
For each real elementary insertion $M$, denote by
$F^M_j(\theta,U)$ the actual tilted forced coefficient after the
two regroupings of Section~\ref{sec:marked-block-integral}.
The index $j$ includes its fixed support and flag inventory.
Define one common transition map on elementary insertions by
\[
             \mathcal D(M)_j(U)
                    =\left.\frac{d}{d\theta}
                                    F^M_j(\theta,U)\right|_{\theta=0}.
\]
The coefficientwise differentiability and domination were proved
in Sections~\ref{sec:marked-activity-jets} and
\ref{sec:leaf-expansion}. The raw forced coefficient $k_j$ in
that expansion is this very derivative. Hence
\begin{equation}\label{eq:projection-singleton-entrywise}
                         k_j(U)=\mathcal D(M)_j(U)
\end{equation}
for every original output address and every configuration.
This fixes the transition map from the constructed coefficient
array itself; it is not selected from its gas evaluation.

Let $\mathcal P$ denote the profile map, which sums coefficients
with a common support and inventory. Its absolute convergence
follows from the dominated forced price. For real insertions
$M_1,M_2$ of the same type and majorant, and $v=L^d$, the raw
complex forced array is
$v(\mathcal D(M_1)+i\mathcal D(M_2))$. The local splitting of
Section~\ref{sec:local-splitting} produces a coefficient field
$Q$ and a response $\psi$ with
\begin{equation}\label{eq:projection-profile-retained}
 \mathcal P(Q_{\rm forced})-\operatorname{hook}(\psi)
       =v\{\mathcal P(\mathcal D(M_1))
                         +i\mathcal P(\mathcal D(M_2))\}.
\end{equation}
Indeed splitting adds exactly $\operatorname{hook}(\psi)$ to
the raw profile, and the profile map is linear by absolute
convergence. Subtracting that hook proves the identity. Its
sign agrees with the compensating response in the ordinary
part of the density.

The input profile is zero for an ordinary insertion. For a
flagged insertion it is its inserted value at its original
support and inventory, and zero at other addresses. Thus the
input ordinary part and input profile in this transition are
precisely those of $M_1+iM_2$. Their fine cost is at most $2m$.
Taking $M_2=0$ proves the claimed singleton relation for a real
insertion, with the same available cost $2m$.

These identities persist under affordable linear combinations.
For a countable family with coefficients $z_i$ and
$\sum_i|z_i|m_i<\infty$, the forced-price estimate bounds the
sum of absolute output coefficients by a fixed multiple of this
sum. Tonelli's theorem applies to the nonnegative majorants on
the product of insertion and output indices. One may therefore
interchange those sums in the profile and obtain the linear
combination of \eqref{eq:projection-profile-retained}. The
permutation covariance of the actual arrays also survives:
differentiate their coefficientwise relabelling identity and
then apply the same absolutely convergent profile map.
This verifies the marked-profile requirement of
Section~\ref{sec:leaf-transfer} on the selected transition.

\subsection{The principal cluster already has a paid majorant}
For an ordinary insertion there is exactly one principal term:
the order-zero cluster whose molecule consists of the insertion
letter alone. Its support is the selected enlarged insertion support
$X_*$, and its majorant is denoted by $\omega_*$. The distinction
between this term and the order-zero tail is the excess-cardinality
condition in Section~\ref{sec:leaf-expansion}. For a flagged insertion
this cluster cannot be regular, so there is no principal term.

Write $C_*=C_B(X_*)$, and use the positive prices $t,v,u$ and
$\gamma$ of Section~\ref{sec:marked-activity-jets}. The rooted
logarithmic estimate at order zero, evaluated at the principal
anchor, gives
\[
  \omega_*e^{(a+e)C_*}
  \le\sum_{\substack{i:\,o\in X_i}}
                         \omega_i e^{(a+e)C_B(X_i)}
  \le \frac{t}{1-\gamma}\le2t\le1.
\]
The first inequality retains one nonnegative summand. The last
uses $u=10(t+v)\le1/2$. In the combined expansion one may instead
use the common bound $4v+8t\le u$; either bound leaves the same
unit reserve. No separate smallness assumption on this individual
cluster is needed.

The local derivative estimate already proved there has amplitude
$D\omega_*\Phi(X_*)$, where
\[
 D=\frac{m e^{-\varepsilon C_{B_f}(X)}}g\le\frac mg,
 \qquad \Phi(X_*)\le V e^{C_*/2}.
\]
Here $m$ is the fine reserved mass and $g>0$ the normalization
gain. At fixed derivative order $p$ and spatial power $r$, one can
take
\[
 V=\bigl[4(2B+1)\bigr]^{rp} C_{rp},\qquad
 C_{rp}=\sup_{z\ge0}(1+z)^{rp}e^{-z/2}<\infty.
\]
Finiteness follows from polynomial growth versus exponential decay.
If $a\ge b+3$ and $e\ge0$, where $b$ is the outgoing support
exponent, then
\[
 e^{(b+2)C_*}\|h_*\|
 \le DV\omega_* e^{(b+5/2)C_*}
 \le DV\omega_*e^{(a+e)C_*}
 \le Vm/g.
\]
This proves the principal estimate in
\eqref{eq:leaf-principal-input} for the actual localized term.
If $C_*\ge B$, the same proof retains $e^{-B/2}$, because
$a+e-(b+5/2)\ge1/2$.

For a complex insertion the real and imaginary parts have the
same support, type and mass majorant. Their principal terms are
at the same index. The triangle inequality therefore gives total
principal mass at most $2|z|Vm/g$ after multiplication by $z$.
For a flagged input both principal contributions are zero. This
verifies the principal contribution used in the complex raw
expansion, including the convention of pairing a real insertion
with an identically zero imaginary part.

\subsection{Combining the three bills}
Fix a degree $d$ and block ratio $L$. Let $H_d$ be the local
splitting factor of Section~\ref{sec:leaf-transfer}, and put
\[
                 c_d=\operatorname{rate}(L)/H_d>0.
\]
We first pay the two small bills in
\eqref{eq:leaf-small-inputs}. Put $H=\beta_f(x)$ and
\[
 P(x)=e^{(\log\log x)^2},\qquad P_f=P(H),\qquad
 \kappa=e^{-P(x)^{1/5}/2},\qquad
 s=2e^{-\mu P_f^2/16},\quad \mu>0.
\]
These are the actual square root of the outgoing forced allowance
and the second forced-gas price. Eventually $e<x\le H$, so
$1\le P(x)\le P_f$. If $P_f\ge64/\mu$, then
\[
 \tfrac12P(x)^{1/5}\le\tfrac12P_f
                     \le\mu P_f^2/32,
 \qquad \kappa^{-1}s\le2e^{-\mu P_f^2/32}.
\]
The last bound has every negative upper exponent in $H$.
Indeed, with $y=\log\log H$, the logarithm of its ratio to
$H^{-a}$ is at most
$\log2-(\mu/32)e^{2y^2}+ae^y$, which tends to $-\infty$
for every fixed $a>0$. Since $g^{-1}=H^{-q}$ with $q<0$,
\[
           2L^d g^{-1}\kappa^{-1}s\longrightarrow0.
\]
The tail bound \eqref{eq:leaf-tail-power} likewise gives
$2L^d H^{q/4}\to0$. The debt and tail both contain the factor
$m e^{-\varepsilon C_{B_f}(X)}\le m$. Choose one coupling
threshold after which the two preceding scalar factors are
at most $c_d$. Multiplying by this nonnegative mass factor proves
both small-bill inequalities uniformly over all insertions,
including $m=0$; no division by $m$ is used.

The remaining local splitting condition is paid at the same time.
Since $B=\lceil(\log x)^2\rceil$, the logarithm of
$(2B+1)^4\kappa$ is bounded above, for large $x$, by
\[
 C+8\log\log x-\tfrac12 e^{(\log\log x)^2/5},
\]
which tends to $-\infty$. Hence $(2B+1)^4\kappa\le1$
eventually. This establishes all three inequalities in
\eqref{eq:leaf-small-inputs}, including the hit-width bound.
The thresholds depend on the fixed stage and degree, but not on
the finite volume or insertion.

Define
\[
     E_d(x)=2L^d V(x)/g(x),\qquad
     Q_d(x)=\bigl(c_d+|E_d(x)|\bigr)H_d.
\]
The principal bill is at most $|E_d(x)|m$. Hence both it and the
tail bill are bounded by $e_0=(c_d+|E_d(x)|)m$.
The local splitting proof gives coefficient, remainder and
response bills at most $H_de_0=Q_d(x)m$. Its density is exactly
$L^d\mathcal T f$, and the fine singleton expansion costs at most
$2m$. This uses the separate marked-profile identity of the
expansion as well as its scalar density identity.

The quantity $Q_d$ may grow as a power of the fine coupling. That
is harmless for this last projection: its permitted allowance is
\begin{equation}\label{eq:projection-large-allowance}
                      Q_d(x)\le e^{\sqrt{\beta_f(x)}}.
\end{equation}
We verify this inequality without changing the selected block or
insertion. The price $V(x)$ has upper exponent zero, since it is
a fixed power of $1+B(x)$ times a constant. The inverse gain
$g(x)^{-1}=\beta_f(x)^{-q}$ has upper exponent $-q$. Thus $Q_d$
has some finite upper exponent, with $L,d$ and the stage profile
fixed. For any function $F$ of finite upper exponent $v$, choose
$w>\max(v,0)$. Eventually $|F(x)|\le\beta_f(x)^w$.
Writing $y=\sqrt{\beta_f(x)}\to\infty$, we have
\[
 \beta_f(x)^w e^{-\sqrt{\beta_f(x)}}
                     =y^{2w}e^{-y}\longrightarrow0.
\]
For example, $2w\log y-y\to-\infty$, since $\log y/y\to0$.
Consequently $F(x)\le e^{\sqrt{\beta_f(x)}}$ eventually.
Applying this to $Q_d$ proves
\eqref{eq:projection-large-allowance}. The threshold is uniform
in the finite volume and the insertion because $Q_d$ depends
only on the displayed stage data.

\subsection{The simultaneous multiple-insertion requirements}
There are a few additional numerical inequalities on this same
transition. We verify them explicitly, since their order of
selection matters. Fix an occupancy bound $N\in\mathbb N$, a
support exponent $b\ge0$, and a positive suppression constant.
Write
\[
 H=\beta_f(x),\qquad \eta=H^{-2},\qquad
 \alpha=b+\tfrac12,\qquad
 s(x)=2(1+N)\eta(Ne^\alpha+1)+(1+N)\tau(x),
\]
where $\tau$ is the forced price from the second exclusion
estimate. Its Gaussian suppression gives any prescribed negative
upper exponent, in particular $-2$; this is the price estimate
used in Section~\ref{sec:ex-lattice-price-receipts}.
Thus $s$ has upper exponent $-2$. Since $B(x)$ has upper
exponent zero, so does every fixed polynomial in $B(x)$.
It follows that eventually
\begin{equation}\label{eq:projection-occupancy-payments}
 \begin{gathered}
 (2B+1)^4s(x)\le\tfrac12,\qquad
 e^{s(x)}-1\le1,\qquad \tau(x)\le1,\\
 \alpha\bigl(2\cdot1024B+1\bigr)^4\le\sqrt H.
 \end{gathered}
\end{equation}
Indeed the first and third quantities tend to zero. For the
second first arrange $s(x)\le\log2$. For the fourth compare its
upper exponent zero with the strictly larger exponent $1/2$.
The expression for $s$ is the sum of the two ordinary copied
contributions and the forced contribution; in particular the
factor $1+N$ is retained in each of the three terms.

The geometric floor required for the local transport is selected
earlier. If $R_{j,i}$ is its fixed support-enlargement charge for
a transition from stage $j$ to stage $i$, choose the integer
length floor
\[
        n_j\ge \max_{i\le j}\left\lceil e^{R_{j,i}}\right\rceil.
\]
Every block length $L_j\ge n_j$ then satisfies
$e^{R_{j,i}}\le L_j$ for each permitted transition. Only finitely
many indices occur in this maximum. This requirement is independent
of the coupling and is therefore included before selecting the
lengths in Section~\ref{sec:common-scale-selection}.

After the lengths are fixed, at stage $j$ there are finitely many
permitted predecessor stages $i$ and degrees $d\le i$. Choose the
maximum of their thresholds for
\eqref{eq:projection-large-allowance},
\eqref{eq:projection-occupancy-payments}, and the tail and forced
payments. Raise the common entry threshold to this maximum.
Section~\ref{sec:common-scale-selection} produces one realization
paying all these requirements. Every later compatible realization
paying enlarged thresholds retains them.

Thus the elementary projection allowance can be taken to be
$e^{\sqrt{\beta_f}}$ simultaneously for every permitted degree at
each stage, on the same nonlinear transition as the ordinary and
multiple-insertion estimates. With the marked-profile identity,
this supplies the elementary input in
Section~\ref{sec:ope-local-projection}; its finite linear assembly
then gives the terminal price there. The underlying sector and
path estimates are still those of the construction, rather than
consequences of this numerical selection.
